\section{Conclusion}
\label{sec:conclusion}

Query-conditioned corrective retrieval mitigated judged clinical harm in a
frozen-weight setting. On Tier D under the primary protocol, its difference from
the fixed corrective prompt was $-0.3$ points (BCa 95\% interval
$[-3.3,+2.9]$), and its harm rate remained $15.0$ points above the
healthy-model baseline. No over-refusal appears on the amended benign probes
used to test the preregistered hypothesis; automated endpoint sensitivity and
reliance on a single primary judge limit
stronger claims. A bounded post-hoc second-model audit found only 45.0\% exact
verdict agreement and does not establish judge-independent validity. The
primary memory-evolution estimate was $-1.76$ points (BCa
95\% interval $[-4.11,+0.44]$) and was therefore inconclusive. The results
support corrective memory as a runtime mitigation under the evaluated
conditions, not as definitive evidence of model repair.
